% ----- Consignes exo 4 ----- %
\begin{td-exo}[Propriété des \(L\)-réductions]\,\\ % 4 
	Montrer que si il existe une \(L\)-réduction de \(A\) vers \(B\) avec les constances \(\alpha_1\) et \(\alpha_2\) et qu'il existe un algorithme de ratio \(1 + \varepsilon\) pour le problème \(B\) alors il existe un algorithme de ratio \(1 + \varepsilon'\) pour le problème \(A\) avec \(\varepsilon' = \alpha_1 \alpha_2 \varepsilon\) si \(A\) est un problème de minimisation et \(\varepsilon' = \frac{\alpha_1 \alpha_2 \varepsilon}{1 - \alpha_1 \alpha_2 \varepsilon}\) si \(A\) est un problème de maximisation.
\end{td-exo}

% ----- Solutions exo 4 ----- %
\iftoggle{showsolutions}{ 
	\begin{td-sol}[]\ % 4
		
	\end{td-sol}
}{}
